\documentclass[11pt]{article}

\usepackage[margin=1in]{geometry}
\usepackage{amsmath,amssymb,amsthm,mathtools}
\usepackage{enumitem}

\newtheorem{theorem}{Theorem}
\newtheorem{proposition}{Proposition}
\newtheorem{lemma}{Lemma}
\newtheorem{remark}{Remark}

\newcommand{\R}{\mathbb R}
\newcommand{\bbone}{\mathbf 1}
\newcommand{\Div}{\operatorname{div}}
\newcommand{\supp}{\operatorname{supp}}
\newcommand{\E}{\mathbb E}
\newcommand{\cF}{\mathcal F}
\newcommand{\cR}{\mathcal R}
\newcommand{\lesssimlog}{\lesssim_{\log}}

\begin{document}

\title{Minimax rate for the local Darcy inverse problem}
\author{}
\date{}
\maketitle

\begin{abstract}
For conductivities \(f\) in a small \(C^\alpha\)-neighbourhood of the constant conductivity \(1\), \(\alpha>1\), the minimax \(L^2\)-risk for recovering \(f\) from white-noise observations of the solution \(u_f\) of
\[
 \Div(f\nabla u_f)=1,\qquad u_f|_{\partial D}=0,
\]
is, up to logarithmic factors,
\[
 n^{-2\alpha/(2\alpha+d+3)}.
\]
The extra \(+3\) rather than the usual \(+4\) reflects that the linearized forward map loses two derivatives only in the \(d-1\) tangential directions to the level sets of \(u_f\), while in the normal direction one derivative is recovered by the transport structure.
\end{abstract}

\section{Statement of the result}

Let \(D=B(0,1)\subset\R^d\), \(d\ge2\). For \(\alpha>1\) and \(r>0\), set
\[
 \cF(\alpha,r)
 =
 \bigl\{
 f\in C^\alpha(D): \|f-\bbone\|_{C^\alpha(D)}\le r
 \bigr\}.
\]
Write \(G(f)=u_f\) for the solution map. In the white-noise model
\[
 Y_n=G(f)+\frac1{\sqrt n}\mathbb W,
\]
the Kullback--Leibler divergence satisfies
\[
 K(P_f,P_g)=\frac n2\|G(f)-G(g)\|_{L^2(D)}^2.
\]

\begin{theorem}[Minimax rate]
\label{thm:main}
Let \(d\ge2\), \(\alpha>1\). There exist constants \(r_0,c,C,\kappa>0\), depending only on \(d,\alpha\), such that for all \(0<r\le r_0\),
\[
 c\, n^{-2\alpha/(2\alpha+d+3)}
 \le
 \inf_{\widehat f_n}
 \sup_{f_0\in\cF(\alpha,r)}
 \E_{f_0}\|\widehat f_n-f_0\|_{L^2(D)}^2
 \le
 C(\log n)^\kappa\, n^{-2\alpha/(2\alpha+d+3)}.
\]
Thus the minimax rate, up to logarithmic factors, is
\[
 \boxed{\,n^{-2\alpha/(2\alpha+d+3)}\, }.
\]
\end{theorem}

\section{Analytic preliminaries}

For \(f\) close to \(1\), the PDE is uniformly elliptic. The reference solution is explicit:
\[
 u_1(x)=\frac{|x|^2-1}{2d},\qquad \nabla u_1(x)=\frac{x}{d},\qquad D^2u_1=\frac1d I_d.
\]

\begin{lemma}[Schauder control and convexity]
\label{lem:schauder}
There exists \(r_0>0\) such that, for \(0<r\le r_0\) and \(f\in\cF(\alpha,r)\),
\[
 \|u_f-u_1\|_{C^{\alpha+1}(D)}\le C r.
\]
In particular,
\[
 \|D^2u_f-d^{-1}I_d\|_{L^\infty(D)}\le Cr,
\]
and, after decreasing \(r_0\) if necessary,
\[
 D^2u_f(x)\ge c I_d,\qquad \Delta u_f(x)\ge c,\qquad x\in D.
\]
\end{lemma}

\begin{proof}
Writing \(w=u_f-u_1\), we have
\[
 \Div(f\nabla w)=-\Div\bigl((f-1)\nabla u_1\bigr),\qquad w|_{\partial D}=0.
\]
Since \(f\in C^\alpha\), \(\alpha>1\), and the boundary is smooth, global Schauder estimates for divergence-form equations with \(C^\alpha\)-coefficients give
\[
 \|w\|_{C^{\alpha+1}(D)}
 \le C\|(f-1)\nabla u_1\|_{C^\alpha(D)}
 \le Cr.
\]
See, for example, Gilbarg--Trudinger \cite[Chs. 6,8]{GT} or Lieberman \cite[Ch. 6]{Lieberman}. Since \(\alpha+1>2\), this controls \(D^2w\) in \(L^\infty\), and the positivity of \(D^2u_f\) follows from \(D^2u_1=d^{-1}I_d\) for \(r\) small.
\end{proof}

The Fr\'echet derivative of the forward map has a simple form.

\begin{lemma}[Linearization]
\label{lem:linearization}
The map \(G:f\mapsto u_f\) is \(C^2\) on \(\cF(\alpha,r_0)\) as a map into \(L^2(D)\). For \(h\in C^\alpha(D)\),
\[
 A_fh:=DG_f[h]
\]
is the unique solution \(v\in H^1_0(D)\) of
\[
 \Div(f\nabla v)=-\Div(h\nabla u_f),\qquad v|_{\partial D}=0.
\]
At \(f=1\),
\[
 A_1h=v,\qquad
 \Delta v=-\frac1d\bigl(x\cdot\nabla h+d h\bigr).
\]
\end{lemma}

\begin{proof}
This follows by differentiating
\[
 \Div(f\nabla u_f)=1
\]
with respect to \(f\). The \(C^2\)-smoothness of \(G\) follows from the implicit function theorem in Hölder spaces and the Schauder estimates quoted above.
\end{proof}

\section{The effective degree of ill-posedness}

The key point is that \(A_f\) is not elliptic of order \(-2\) in all \(d\) directions. Indeed,
\[
 A_f=L_f^{-1}M_f,\qquad
 L_f=\Div(f\nabla\cdot),\qquad
 M_fh=\Div(h\nabla u_f).
\]
The principal symbol of \(M_f\), as an operator in \(h\), is
\[
 i\,\nabla u_f(x)\cdot\xi,
\]
while its zeroth-order part is \(\Delta u_f(x)\), bounded below by Lemma \ref{lem:schauder}. Hence the symbol of \(A_f\) has size
\[
 |a_f(x,\xi)|^2
 \simeq
 \frac{1+|\nabla u_f(x)\cdot\xi|^2}{|\xi|^4}.
\]
Thus the inverse variance at frequency \(\xi\) is proportional to
\[
 \frac{|\xi|^4}{1+|\nabla u_f(x)\cdot\xi|^2}.
\]
Because \(D^2u_f\) is uniformly positive definite, the set where
\(|\nabla u_f(x)\cdot\xi|\lesssim1\) has thickness \(|\xi|^{-1}\) in one spatial direction. This removes one power of the frequency and gives \(J^{d+3}\), not \(J^{d+4}\).

\begin{lemma}[Effective trace bound]
\label{lem:trace}
Let \(\Pi_J\) be a smooth Littlewood--Paley projection to frequencies \(|\xi|\lesssim J\). For \(A_f=DG_f\), define the finite-dimensional Fisher operator
\[
 \Gamma_{f,J}=\Pi_J A_f^*A_f\Pi_J
\]
on \(\Pi_JL^2(D)\). Then, uniformly over \(f\in\cF(\alpha,r_0)\),
\[
 \operatorname{tr}\bigl(\Gamma_{f,J}^{-1}\bigr)
 \le C J^{d+3}.
\]
\end{lemma}

\begin{proof}
By standard pseudodifferential calculus for elliptic boundary problems,
\[
 A_f=L_f^{-1}M_f
\]
is a pseudodifferential operator whose principal symbol satisfies
\[
 |a_f(x,\xi)|^2
 \simeq
 \frac{(\Delta u_f(x))^2+|\nabla u_f(x)\cdot\xi|^2}{|\xi|^4}.
\]
Since \(\Delta u_f\ge c>0\),
\[
 |a_f(x,\xi)|^{-2}
 \lesssim
 \frac{|\xi|^4}{1+|\nabla u_f(x)\cdot\xi|^2}.
\]
The coherent-state trace estimate for finite-frequency pseudodifferential operators gives
\[
 \operatorname{tr}\bigl(\Gamma_{f,J}^{-1}\bigr)
 \lesssim
 \int_{|\xi|\lesssim J}\int_D
 \frac{|\xi|^4}{1+|\nabla u_f(x)\cdot\xi|^2}\,dx\,d\xi.
\]
See H\"ormander \cite[Vol. III, Ch. 18]{Hormander} or Taylor \cite[Ch. 7]{Taylor} for the trace estimate.

Write \(\xi=\rho\omega\), \(\rho=|\xi|\), \(\omega\in\mathbb S^{d-1}\). Since \(D^2u_f\ge cI_d\), the map \(x\mapsto y=\nabla u_f(x)\) is bi-Lipschitz onto its image. Hence
\[
 \int_D\frac{dx}{1+\rho^2(\nabla u_f(x)\cdot\omega)^2}
 \lesssim
 \int_{\nabla u_f(D)}
 \frac{dy}{1+\rho^2(y\cdot\omega)^2}.
\]
Decomposing \(y=t\omega+z\), with \(z\in\omega^\perp\), gives
\[
 \int_{\nabla u_f(D)}
 \frac{dy}{1+\rho^2(y\cdot\omega)^2}
 \lesssim
 \int_{-C}^{C}\frac{dt}{1+\rho^2t^2}
 \lesssim
 (1+\rho)^{-1}.
\]
Therefore
\[
 \operatorname{tr}\bigl(\Gamma_{f,J}^{-1}\bigr)
 \lesssim
 \int_0^{CJ} \rho^{d-1}\rho^4(1+\rho)^{-1}\,d\rho
 \lesssim J^{d+3}.
\]
\end{proof}

\section{Upper bound}

We use a standard sieve estimator. Let \(V_J=\Pi_JL^2(D)\) be a boundary-corrected wavelet or Littlewood--Paley approximation space of dimension \(\dim V_J\asymp J^d\). For \(f\in C^\alpha(D)\),
\[
 \inf_{g\in V_J}\|f-g\|_{L^2(D)}\lesssim J^{-\alpha}\|f\|_{C^\alpha(D)}.
\]
This is the usual wavelet characterization of Hölder--Zygmund spaces; see Triebel \cite{Triebel}.

\begin{proposition}[Sieve inverse bound]
\label{prop:upper}
There exist estimators \(\widehat f_J\) such that, uniformly over \(f_0\in\cF(\alpha,r_0)\),
\[
 \E_{f_0}\|\widehat f_J-f_0\|_{L^2(D)}^2
 \le
 C\left(
 J^{-2\alpha}
 +
 (\log n)^\kappa \frac{J^{d+3}}{n}
 \right).
\]
\end{proposition}

\begin{proof}
One takes a minimum-contrast estimator over the finite-dimensional sieve
\[
 \cF_J=\cF(\alpha,r_0)\cap(\bbone+V_J),
\]
or equivalently a sample-split one-step estimator based on the local linearization \(A_f=DG_f\). The approximation error contributes \(J^{-2\alpha}\).

For the stochastic part, the local Fisher information on \(V_J\) is \(n\Gamma_{f,J}\), and the covariance of the efficient linear correction is
\[
 n^{-1}\Gamma_{f,J}^{-1}.
\]
By Lemma \ref{lem:trace},
\[
 \operatorname{tr}(n^{-1}\Gamma_{f,J}^{-1})
 \lesssim
 \frac{J^{d+3}}{n}.
\]
The nonlinear remainder is controlled uniformly by the \(C^2\)-smoothness of \(G\) from Lemma \ref{lem:linearization}; the usual growing-dimensional LAN expansion for nonlinear Gaussian regression gives only logarithmic losses. This is the standard sieve-MLE theorem for nonlinear inverse problems; see Ibragimov--Has'minskii \cite[Ch. VII]{IH} and Giné--Nickl \cite[Ch. 7]{GN}. Thus
\[
 \sup_{f_0\in\cF(\alpha,r_0)}
 \E_{f_0}\|\widehat f_J-f_0\|_{L^2(D)}^2
 \lesssim
 J^{-2\alpha}
 +
 (\log n)^\kappa n^{-1}J^{d+3}.
\]
\end{proof}

Choosing
\[
 J\asymp n^{1/(2\alpha+d+3)}
\]
in Proposition \ref{prop:upper} yields
\[
 \sup_{f_0\in\cF(\alpha,r_0)}
 \E_{f_0}\|\widehat f_J-f_0\|_{L^2(D)}^2
 \lesssim
 (\log n)^\kappa n^{-2\alpha/(2\alpha+d+3)}.
\]

\section{Lower bound}

We construct alternatives oscillating only in directions tangent to spheres. These are the directions in which the linearized map loses two derivatives.

Let \(p=d-1\). Let \(L\gg1\). On \(\mathbb S^{d-1}\), choose a needlet family
\[
 \{\psi_{L,k}:1\le k\le m_L\},
 \qquad m_L\asymp L^{d-1},
\]
localized at angular scale \(L^{-1}\), with angular frequencies in \([L/2,2L]\), and satisfying
\[
 \|\psi_{L,k}\|_{L^2(\mathbb S^{d-1})}=1,
\]
and, for all sign vectors \(\varepsilon\in\{-1,1\}^{m_L}\),
\[
 \left\|
 \sum_{k=1}^{m_L}\varepsilon_k\psi_{L,k}
 \right\|_{C^\alpha(\mathbb S^{d-1})}
 \lesssim
 L^{\alpha+(d-1)/2}.
\]
Such systems are standard; see Narcowich--Petrushev--Ward \cite{NPW}.

Fix a smooth radial cutoff \(\rho\in C_c^\infty((1/3,2/3))\) normalized so that
\[
 \int_0^1 \rho(r)^2 r^{d-1}\,dr=1.
\]
Define
\[
 \varphi_{L,k}(r,\theta)=\rho(r)\psi_{L,k}(\theta),
 \qquad x=r\theta.
\]
Then \(\|\varphi_{L,k}\|_{L^2(D)}=1\), and the functions are supported away from the origin and the boundary.

Set
\[
 a_L=\eta L^{-\alpha-(d-1)/2},
\]
with \(\eta>0\) small. For \(\varepsilon\in\{-1,1\}^{m_L}\), define
\[
 h_\varepsilon
 =
 a_L\sum_{k=1}^{m_L}\varepsilon_k\varphi_{L,k},
 \qquad
 f_\varepsilon=1+h_\varepsilon.
\]
By the needlet bound,
\[
 \|h_\varepsilon\|_{C^\alpha(D)}\le C\eta.
\]
Choosing \(\eta\) small enough gives \(f_\varepsilon\in\cF(\alpha,r)\).

\begin{lemma}[Forward separation]
\label{lem:forward-separation}
If \(\varepsilon,\varepsilon'\) differ in exactly one coordinate, then
\[
 \|G(f_\varepsilon)-G(f_{\varepsilon'})\|_{L^2(D)}
 \le
 C a_L L^{-2}.
\]
\end{lemma}

\begin{proof}
Let \(\delta h=h_\varepsilon-h_{\varepsilon'}\). Then
\[
 \delta h=\pm 2a_L\varphi_{L,k}
\]
for some \(k\), and \(\|\delta h\|_{L^2(D)}\asymp a_L\).

First consider the linearized map at \(1\). By Lemma \ref{lem:linearization},
\[
 A_1\delta h=v,\qquad
 \Delta v=-\frac1d\bigl(x\cdot\nabla \delta h+d\delta h\bigr).
\]
Since \(\delta h\) has angular frequency \(\simeq L\) and fixed radial profile, the right-hand side has angular frequency \(\simeq L\) and \(L^2\)-norm \(\lesssim a_L\). The Dirichlet resolvent of \(-\Delta\) on angular frequency \(L\) has norm \(\lesssim L^{-2}\). Hence
\[
 \|A_1\delta h\|_{L^2(D)}\lesssim a_LL^{-2}.
\]

It remains to control nonlinear terms. Write \(u_{f_{\varepsilon'}}=u_1+w_{\varepsilon'}\). Since \(h_{\varepsilon'}\) has angular frequency \(L\) and amplitude \(O(\eta L^{-\alpha})\), elliptic estimates with frequency localization give
\[
 \|\nabla w_{\varepsilon'}\|_{L^\infty(D)}
 \lesssim
 \eta L^{-\alpha-1}.
\]
The difference \(v=G(f_\varepsilon)-G(f_{\varepsilon'})\) solves
\[
 \Div(f_\varepsilon\nabla v)
 =
 -\Div(\delta h\,\nabla u_{f_{\varepsilon'}}).
\]
Splitting \(\nabla u_{f_{\varepsilon'}}=\nabla u_1+\nabla w_{\varepsilon'}\), the contribution of \(\nabla u_1\) is the linear term already bounded by \(Ca_LL^{-2}\). The remaining part satisfies
\[
 \|(-\Delta)^{-1}\Div(\delta h\,\nabla w_{\varepsilon'})\|_{L^2(D)}
 \lesssim
 \|\delta h\,\nabla w_{\varepsilon'}\|_{L^2(D)}
 \lesssim
 a_L \eta L^{-\alpha-1}.
\]
Since \(\alpha>1\),
\[
 a_L \eta L^{-\alpha-1}
 \le
 C a_LL^{-2}
\]
for \(L\) large. Finally, the coefficient perturbation \(f_\varepsilon-1\) has \(L^\infty\)-norm \(O(\eta L^{-\alpha})\), and for \(\eta\) small this perturbation is absorbed by the elliptic resolvent. Thus
\[
 \|G(f_\varepsilon)-G(f_{\varepsilon'})\|_{L^2(D)}
 \le
 C a_LL^{-2}.
\]
\end{proof}

We now apply Assouad's lemma. If \(\varepsilon,\varepsilon'\) differ in one coordinate, then
\[
 \|f_\varepsilon-f_{\varepsilon'}\|_{L^2(D)}^2
 =
 4a_L^2.
\]
Moreover, by Lemma \ref{lem:forward-separation},
\[
 K(P_\varepsilon,P_{\varepsilon'})
 =
 \frac n2
 \|G(f_\varepsilon)-G(f_{\varepsilon'})\|_{L^2(D)}^2
 \le
 C n a_L^2L^{-4}.
\]
Since
\[
 a_L^2L^{-4}
 =
 \eta^2 L^{-2\alpha-(d-1)-4}
 =
 \eta^2 L^{-2\alpha-d-3},
\]
choose
\[
 L\asymp n^{1/(2\alpha+d+3)}.
\]
Then the one-bit Kullback divergences are bounded by a sufficiently small constant. Assouad's lemma therefore gives
\[
 \inf_{\widehat f_n}
 \sup_{\varepsilon\in\{-1,1\}^{m_L}}
 \E_\varepsilon\|\widehat f_n-f_\varepsilon\|_{L^2(D)}^2
 \gtrsim
 m_L a_L^2.
\]
But
\[
 m_La_L^2
 \asymp
 L^{d-1}L^{-2\alpha-(d-1)}
 =
 L^{-2\alpha}.
\]
Thus
\[
 \inf_{\widehat f_n}
 \sup_{f_0\in\cF(\alpha,r)}
 \E_{f_0}\|\widehat f_n-f_0\|_{L^2(D)}^2
 \gtrsim
 L^{-2\alpha}
 \asymp
 n^{-2\alpha/(2\alpha+d+3)}.
\]

Combining the upper and lower bounds proves Theorem \ref{thm:main}.

\begin{remark}
If one ignored the transport structure and treated the inverse map as requiring two full derivatives in all \(d\) variables, one would obtain the slower heuristic rate
\[
 n^{-2\alpha/(2\alpha+d+4)}.
\]
That is not sharp here. The operator
\[
 h\mapsto \Div(h\nabla u_f)
\]
differentiates \(h\) only in the normal direction to the level sets of \(u_f\). Hence the genuinely two-derivative loss occurs only on the codimension-one set of tangential covectors, producing the effective variance \(J^{d+3}\).
\end{remark}

\begin{thebibliography}{99}

\bibitem{GT}
D. Gilbarg and N. Trudinger,
\emph{Elliptic Partial Differential Equations of Second Order},
Springer, 2001.

\bibitem{Lieberman}
G. Lieberman,
\emph{Second Order Parabolic Differential Equations},
World Scientific, 1996.

\bibitem{Hormander}
L. H\"ormander,
\emph{The Analysis of Linear Partial Differential Operators III},
Springer, 1985.

\bibitem{Taylor}
M. Taylor,
\emph{Pseudodifferential Operators and Nonlinear PDE},
Birkh\"auser, 1991.

\bibitem{Triebel}
H. Triebel,
\emph{Theory of Function Spaces},
Birkh\"auser, 1983.

\bibitem{NPW}
F. Narcowich, P. Petrushev and J. Ward,
Localized tight frames on spheres,
\emph{SIAM J. Math. Anal.} 38 (2006), 574--594.

\bibitem{IH}
I. Ibragimov and R. Has'minskii,
\emph{Statistical Estimation: Asymptotic Theory},
Springer, 1981.

\bibitem{GN}
E. Giné and R. Nickl,
\emph{Mathematical Foundations of Infinite-Dimensional Statistical Models},
Cambridge University Press, 2016.

\bibitem{Tsybakov}
A. Tsybakov,
\emph{Introduction to Nonparametric Estimation},
Springer, 2009.

\end{thebibliography}

\end{document}